\begin{explanationbox}
	\textbf{\textcolor{CExplanation}{一句话直觉}}

	\medskip
	这个公式用一个分式快速算出一条直线与一个平面的交点比例参数，通过参数范围判断是否真正相交于棱上。

	\medskip
	\textbf{\textcolor{CExplanation}{核心拆解}}
	\begin{description}[style=nextline, leftmargin=2.8em, labelsep=0.5em, itemsep=0.45em]
		\item[\textbf{要点一（交参分式结构）：}]
			分子 $\vec{n}\cdot(\vec{p}-\vec{a})$ 是点 $A$ 到平面 $\pi$ 的带号距离（以法向量为度量），分母 $\vec{n}\cdot\vec{v}$ 是方向向量在法向量上的投影。两商即得直线穿过平面的比例。
		\item[\textbf{要点二（范围确定截取）：}]
			$t\in[0,1]$ 表示交点落在从 $A$ 到 $B$ 的线段内，对应棱上实交点；否则交点在棱的延长线上，截面与棱不相交。
	\end{description}

	\medskip
	\textbf{\textcolor{CExplanation}{几何本质（平面截距法）}}

	\medskip
	将直线参数方程 $\vec{x}=\vec{a}+t\vec{v}$ 代入平面方程 $\vec{n}\cdot(\vec{x}-\vec{p})=0$，解出唯一参数 $t$，本质是直线穿过平面的瞬间位置。

	\medskip
	\textbf{\textcolor{CExplanation}{代数意义（比例因子）}}

	\medskip
	$t$ 表示从起点 $A$ 到交点的位移占全位移 $\vec{v}$ 的比例，是纯数量，只依赖向量的点积，无需解方程组。

	\medskip
	\textbf{\textcolor{CExplanation}{考点价值（截面作图计算）}}
	\begin{itemize}[leftmargin=*, itemsep=0.5em, topsep=0.4em]
		\item \textbf{考法一（截面顶点定位）：} 给出几何体和截面三点，直接使用公式求各棱交点，进而得到截面多边形。
		\item \textbf{考法二（空间存在性判断）：} 利用 $t\notin[0,1]$ 判断平面与棱不相交，排除非顶点。
	\end{itemize}

	\medskip
	\textbf{\textcolor{CExplanation}{顿悟点}}

	\medskip
	\textbf{把截面作图转化为代数计算：} 不再需要手动画图判断交点在棱上还是延长线上，公式给出精确比例。

	\medskip
	\textbf{\textcolor{CExplanation}{使用场景}}
	\begin{itemize}[leftmargin=*, itemsep=0.5em, topsep=0.4em]
		\item \textbf{场景一（计算截面积）：} 先算各棱交点，再按多边形面积公式算截面面积。
		\item \textbf{场景二（判断截面范围）：} 空间中平面切割几何体，通过 $t$ 范围确定哪些棱被切割。
	\end{itemize}

\end{explanationbox}
